% ============================================
%  preface.tex — 译者序
% ============================================
\chapter*{第三版前言}
\addcontentsline{toc}{chapter}{第三版前言}

第二版以来的主要变化，是本版加入了 Gröbner 基。其基本理论在新的第 9.6 节中介绍。求解多项式方程组的应用（消元理论）出现在该节末尾，从而使这一节成为关于此主题的一个自足基础。随后，更多的应用和例子被融入第 15 章对仿射代数集和 \(k\)-代数同态的处理中。虽然后一章中的理论仍然不依赖于 Gröbner 基，但新增的应用、例子和计算技巧显著加强了该部分的展开；我们建议把第 9.6 节作为进入第 15 章的过渡来阅读，或者与第 15 章并行阅读。第 9.6 节和第 15 章中加入了大量涉及 Gröbner 基的习题，这些习题既有计算性质的，也有理论性质的。关于 Gröbner 基的初步习题可以（而且为了帮助理解算法，应当）手算完成；但较大的计算，特别是第 15 章习题中 Gröbner 基的大部分使用，很可能需要借助计算机完成。

其他变化包括对 168 阶单群分类（第 6.2 节）的处理作了简化，并加入了一个通过 2 阶射影平面给出的唯一性证明。文本中的其他若干证明或部分内容也作了轻微修订。全书各处加入了许多新的习题，主要放在各节末尾，以便尽可能保留早期版本的编号方案。特别地，新增习题涉及非交换环上的自由模（10.3）、Krull 维数（15.3）以及平坦模（10.5 和 17.1）。

与以前各版一样，本书所包含的内容远多于通常一年课程所能讲授的内容。一门基本的入门（一年）课程大概应包括第一部分直到第 5.3 节、第二部分直到第 9.5 节、第 10.1、10.2、10.3、11.1、11.2 节以及第四部分。第 12 章也应讲授，可以放在第四部分之前或之后。第 5、6、9、10 和 11 章中的其他主题可以穿插在这样的课程中，也可以在课程末尾视时间许可讲授。

第 10.4 和 10.5 节的难度比第 10 章开头几节略高；对于按顺序阅读本书的读者，初读时可以推迟阅读。后一节关于正合列性质的内容虽然相当长，但通过在各个小节中对三个基本函子的平行处理而保持了连贯性。

在核心材料之外，第三版为希望学习现代代数若干重要领域的学生和教师提供了相当大的灵活性，无论是作为独立学习还是作为课程内容。例如，对于已有一定代数背景的学生，可以组织结构良好的一学期课程，内容包括：第 9.6 节以及第 15、16 章；或第 10、17 章；或第 5、6 章以及第六部分。这些选择中的任何一种，也都会为后续课程打下坚实基础，使学生能够更深入地学习许多可能领域中的某一个，例如代数数论、代数拓扑、代数几何、表示论、Lie 群等等。

新增材料的选择，以及将这些材料展开并融入正文的方式，与本书的一个基本主题相一致：数学不同领域之间丰富的相互作用所产生的力量与美。全书始终强调尽可能用许多例子来激发重要代数概念的引入和发展。我们并未试图做到百科全书式的包罗万象，而是试图以一种显示这些思想十分自然地发展出来的方式，触及初等代数中的许多核心主题。

许多重要的思想和结果出现在习题中。这并不是因为它们不重要，而是因为它们不容易放入正文的行文脉络中，但又太重要而不能完全略去。关于同一主题的一系列习题前面会有一些说明，并且经过安排，使读者即使不实际做这些习题也可以阅读。在某些情况下，新材料首先在习题中引入——往往比正文中出现早几节——这样学生可以通过做这些习题而较容易地初步接触这些材料（例如，Lagrange 定理出现在第 1.7 节的习题中，而在正文中出现在第 3.2 节）。所有习题都在正文范围之内；在我们认为需要的地方，提示用[方括号]给出。我们认为可能不那么直接的习题，通常会以给出习题答案的方式来表述；同时，许多习题被分解成一串较为常规的习题，以便使其更容易处理。

我们也有意尽量减少正文中的函子语言，以便使叙述尽可能保持初等性。对于进一步阅读的资料，在已有许多优秀选择可供参考的地方，我们没有给出具体文献。并且，虽然我们已经尽力包含尽可能多的基本主题，但若由于篇幅或个人取舍的原因而遗漏了读者特别喜爱的某些内容，我们在此致歉。

我们深切感谢，并愿在这里向世界各地的许多学生和同事表示感谢；十五年多以来，他们提供了宝贵的意见、见解和鼓励——他们持续的关注激励我们写成了这个第三版。

\begin{flushright}
	David Dummit\\
	Richard Foote\\
	June, 2003
\end{flushright}
\chapter*{译者序}
\addcontentsline{toc}{chapter}{译者序}

本文是 David S.\ Dummit 与 Richard M.\ Foote 所著
\textit{Abstract Algebra}（第三版）的中文翻译稿。

该书是抽象代数领域的经典教材，涵盖群论、环论、模论、
域论与 Galois 理论、交换代数、同调代数以及有限群表示论
等核心内容。全书共 19 章，分为六大部分，体系完整，习题丰富。

翻译过程中力求忠实于原著，同时兼顾中文数学表达的流畅性。


\begin{flushright}
	万一之 \\
	\today
\end{flushright}